\section{Cache}
Currently, the dominant memory technology in the PC market is the Double Data
Rate-Synchronous Dynamic RAM, delivering memory on both the positive and
negative edge of the clock.
\begin{example}[DDR5]
  The DDR5 Generation delivers information at a rate of
  \[MemClkFreq \times 2^{5} \times 8 \text{ words}.\] 
\end{example}
The technologies of our processors and ICs have been increasing exponentially,
introducing a gap with the performance of the DRAMs. This necessitates the
utilisation of faster memory technology such as SRAM (with lower density and
higher costs, but faster access latency).
\begin{figure}[h!]
  \centering
  \includegraphics[width=0.5\textwidth]{memories2}
  \includegraphics[width=0.4\textwidth]{memories}
\end{figure}

To simulate a large and fast access memory, we use a hierachy of memory
technologies,to keep frequently and recently used data in a smaller but fast
memory and refer to the larger and slower memory only if needed.

\paragraph{Principle of Locality} This is effective as programs access only a
small portion of the memory address space within a small time interval.
\begin{figure}[h!]
  \centering
  \includegraphics[width=0.3\textwidth]{workingSet}
\end{figure}
\begin{itemize}
  \item \textbf{Temporal Locality}: When an item is referenced, it tends to be
    referenced again. This applies for many programs involving loops.
  \item \textbf{Spatial Locality}: When an item is referenced, nearby items
    tend to be referenced as well. This especially applies to certain data
    structures, particularly arrays and instructions.
\end{itemize}

\begin{definition}[Working Set]
  A working set is a set of locations accessed during some $\Delta t$.
\end{definition}
We may use different working sets during different phases of execution, and we
aim to capture the working set to keep in memory with fastest access latencies.
We do so by creating a cache that is hardware managed.

We now investigate various implementations of mappings betweeen the memory and
the cache.

\begin{definition}[Cache Blocks/Lines]
  A \textit{cache block} is a unit of transfer between the memory and cache,
  typically consisting more than one word, to take advantage of spatial locality.
\end{definition}
\begin{figure}[h!]
  \centering
  \includegraphics[width=0.6\textwidth]{cache_block}
\end{figure}

\subsection{Cache Misses}
A cache miss occurs due to one of three reasons:
\begin{itemize}
  \item \textbf{Compulsory Miss} (Cold Start Miss / First Reference Miss)
    occurs on first access to the block. The block must be brought into the
    cache.
  \item \textbf{Conflict Miss} (Collision Miss / Interference Miss) occurs in
    the case of directed mapped cache or set associative cache, when several
    blocks are mapped to the same block or set, and it is not a compulsory
    miss.
  \item \textbf{Capacity Miss} occurs when blocks were discarded as the cache
    cannot contain all blocks requested.
\end{itemize}

\subsection{Memory Access}
When a memory address is specified, the data is either in the cache, or not.
\begin{definition}[Cache Hit]
  A \textit{cache hit} occurs when the data is in the cache. We call the fraction
  of memory accesses that hit the \textit{hit rate}, and the time to access
  the cache the \textit{hit time}.
\end{definition}
\begin{definition}[Cache Miss]
  A \textit{cache miss} occurs when the data is not in the cache. We call
  the fraction of memory accesses that miss the \textit{miss rate}$=1-$ hit
  rate, and the time to replace the cache block and then to access the cache
  the  \textit{miss penalty}.
\end{definition}

The average access time is hence obtained by
\[Average\ Access\ Time= Hit\ Rate \times Hit\ Time + Miss\ Rate \times Miss\
Penalty.\] 

\paragraph{Block Size Trade-off}
\begin{figure}[h!]
  \centering
  \includegraphics[width=0.6\textwidth]{blockTrade}
\end{figure}
With a larger block size:
\begin{itemize}
  \item Take advantage of spatial locality
  \item Larger miss penalty, requiring longer time to fill up the block
  \item Block size too large relative to cache size results in too few cache
    blocks, increasing miss rate
\end{itemize}

\subsubsection{Reading Data}
On reading data, we first check the cache for the existence of the data block.
If there is a cache miss, we request the data from the main memory, write into
the cache, and then load from the cache to the register.
\begin{figure}[h!]
  \centering
  \includegraphics[width=0.5\textwidth]{cacheRead}
\end{figure}

\subsubsection{Writing Data}
On writing data, in case of a \textbf{cache hit}, we must maintain consistency
between the cache and main memory. As such, we have the write policies:
\begin{enumerate}
  \item \textbf{Write-Through Cache}: Write data to both the cache and main
    memory. We add a write buffer between the cache and main-memory, that is
    operated by a memory controller to allow for processor to operate at the
    speed of the cache and write buffer instead of the main memory.
    \begin{figure}[h!]
      \centering
      \includegraphics[width=0.4\textwidth]{write_through}
    \end{figure}
    \vspace{-4mm}
  \item \textbf{Write-Back Cache}: Write data to only cache. Commit to memory
    only when the cache block is replaced. We add an additional (Dirty) bit to
    each cache block, to indicate if the cache block has been modified. We only
    attempt to write back to memory when the cache block is being evicted, and
    its dirty bit is 1.
\end{enumerate}

In case of a \textbf{cache miss}, we have the write miss policies:
\begin{enumerate}
  \item \textbf{Write Allocate}: Load the complete block into the cache, change
    only the required word in the cache, and write to main memory depending on
    write policy.
  \item \textbf{Write Around}: Write directly to main memory only.
\end{enumerate}

\begin{figure}[h!]
  \centering
  \includegraphics[width=0.5\textwidth]{cacheWrite}
\end{figure}

\subsection{Direct Mapped Cache}
In the Direct Mapped Cache, we identify the next LSBs after
the byte offsets as the \textit{cache index}, which identifies the line in the
cache where this block can be stored. Since multiple blocks in memory can have
the same cache index, we further specific the address with the \textit{tag},
the rest of the bits in the memory address and we have a mapping from the main
memory to the cache.
\[Cache\ Index=(Block\ Number) \bmod (Number\ of\ Cache\ Blocks).\] 
\[Tag=(Block\ Number) / (Number\ of\ Cache\ Blocks).\] 
\begin{figure}[h!]
  \centering
  \includegraphics[width=0.55\textwidth]{dmc_map}
\end{figure}\\
The cache hence consists of the following components for each block:
\begin{itemize}
  \item \textbf{Valid Bit} indicating whether the cache line contains valid data
  \item \textbf{Tag} of the memory block
  \item \textbf{Data Block}
\end{itemize}
Given a memory address, we compute its index and tag, and a cache hit occurs $\iff$
\verb|Valid[index]=TRUE| $\wedge$ \verb|Tag[index]=tag|.
\begin{figure}[h!]
  \centering
  \includegraphics[width=0.5\textwidth]{cacheCircuit}
\end{figure}

\subsection{Set Associative Cache}
In a Set Associative Cache, we attempt to circumvent the issue of conflict
misses. We define a $N$-way Set Associative Cache in which each memory block
can be placed in a fixed number ($N>1$) of locations in the cache. The cache
consists a number of sets, each containing $N$ cache blocks. Each memory block
maps to a unique cache set within which it can be placed in any of the $N$
cache blocks.

\begin{figure}[h!]
  \centering
  \includegraphics[width=0.6\textwidth]{sacMap}
\end{figure}

\begin{figure}[h!]
  \centering
  \includegraphics[width=0.65\textwidth]{sacCircuit}
\end{figure}

\begin{prop}[Miss Rate of SAC vs DAC]
  A Direct-Mapped Cache of size $N$ has about the same miss rate as a $2$-way
  Set Associative Cache of size $\frac{N}{2}$.
\end{prop}

\subsubsection{Block Replacement Policy}
When the cache set is full, we have a choice of memory blocks that can be replaced
within the cache set. As such, we may define a block replacement policy:
\begin{itemize}
  \item \textbf{Least Recently Used} (Difficult to keep track if $N$ is large)
  \item \textbf{First In First Out}
  \item \textbf{Random Replacement}
  \item \textbf{Least Frequently Used}
\end{itemize}

\subsection{Fully Associative Cache}
In a Fully Associative Cache, a memory block can be placed in any location in
the cache. This eliminates conflict misses, since blocks are not restricted in
their locations. However, this requires searching all cache blocks for memory
access. Essentially, the entire cache is a cache set for all memory blocks.

\begin{figure}[h!]
  \centering
  \includegraphics[width=0.55\textwidth]{facmap}
\end{figure}
\begin{figure}[h!]
  \centering
  \includegraphics[width=0.6\textwidth]{facCircuit}
\end{figure}

\subsection{Cache Performances Comparisons}
In comparison of cache sizes and associativity, we may derive various trends
that allow us to optimise caches for specific contexts.
\paragraph{Compulsory Misses} remain the same regardless of cache size and
associativity.
\paragraph{Conflict Misses} decreases with increasing associativity, for the
same cache size. It is $0$ for FACs.
\paragraph{Capacity Misses} remain the same regardless of associativity, for
the same cache size. It decreases with increasing cache size.
\[Total\ Miss=Cold\ Miss + Conflict\ Miss + Capacity\ Miss.\] 

\subsection{Multilevel Caches}
To decrease average access times, we may also decrease the miss penalty. One
such method is to consider multilevel caches, to either separate data and
instruction caches (with differing localities) or a unified cache.

\begin{figure}[h!]
  \centering
  \includegraphics[width=0.7\textwidth]{multilevel}
\end{figure}
